\section{The Angels in the Liturgy}\label{sec:liturgy}

The Church prays what she believes about the angels. At every Mass she
sings with them, names Michael among those whose prayer she begs, blesses
her incense by his intercession, and asks that her sacrifice be borne to
the altar on high by the hand of an angel. In the Office she keeps their
feasts with the lessons of Gregory, Jerome, Bernard, and Hilary; at night
she asks that the holy angels dwell in her houses; at death she sends the
soul forth in the name of the angelic orders and hands it to Michael, the
standard-bearer who presents souls in the holy light. What Scripture
teaches and the Fathers expound about the angels' praise, their ministry
to men, and their office at death, the Roman liturgy turns into
prayer.\footnote{The Roman books cited are the Missal of the typical
edition of 1920 (\emph{MR}~1920), in the Vatican printing and in Mame's
Tours printing \latin{iuxta typicam} of 1922, whose pages are given; the
Missal of 1962 (\emph{MR}~1962); Cummiskey's Philadelphia lay missal of
1861, which sets the Latin of the Order of Mass beside an English
translation; the Roman Breviary printed at Tours in 1898 (\emph{BR}), with
the English of John, Marquess of Bute (1908); the \emph{Rituale Romanum}
of Ratisbon (1872, \emph{RR}), with the English of the \emph{Manual of
Prayers} published by order of the Third Plenary Council of Baltimore
(1889); and, for the postconciliar calendar and readings, the Missal of
Paul~VI (third typical edition, 2002, emended 2008) with the \emph{Ordo
lectionum Missae} of 1981.} The Byzantine liturgy of the angels belongs to
Section~\ref{sec:byzantine}; the dates and ranks of every feast stand
together in Appendix~\ref{app:calendar}.

\sectionguard
\subsection{The Mass Sung with the Angels}

\paragraph{The hymn of the angels.} The first words of the \latin{Gloria}
are the angels' own. On the night of the Nativity ``suddenly there was
with the angel a multitude of the heavenly army, praising God and saying:
Glory to God in the highest: and on earth peace to men of good will''
(Luke 2:13--14). The Church takes up their song at the altar,
\latin{Gloria in excelsis Deo, et in terra pax hominibus bonae voluntatis},
``Glory be to God on high, and on earth, peace to men of good will''
(Cummiskey, \emph{Roman Missal} 1861, p.~xix), and continues it in her own
praise of the Father and of the Lamb. Aquinas reads the hymn as the part of
the Mass that ``commemorates the heavenly glory, to the possession of
which, after this life of misery, we are tending''; it is therefore sung
``on festival days, on which the heavenly glory is commemorated, but [is]
omitted in those sorrowful offices which commemorate our unhappy state''
(\ST{III q.~83 a.~4 co.}). The angels' hymn of the Nativity is the
Church's hymn of the heavenly homeland.

\paragraph{Michael in the Confiteor.} Before she goes up to the altar the
Church confesses her sins to God and to the court of heaven, and the first
named after the Mother of God is the Archangel:
\begin{quote}
\latin{Confiteor Deo omnipotenti, beatae Mariae semper Virgini, beato
Michaeli Archangelo \dots\ Ideo precor beatam Mariam semper Virginem,
beatum Michaelem Archangelum \dots\ orare pro me ad Dominum Deum nostrum.}
\end{quote}
``I confess to Almighty God, to blessed Mary, ever a Virgin, to blessed
Michael the Archangel, to blessed John Baptist, to the holy Apostles Peter
and Paul, to all the saints \dots\ Therefore I beseech the Blessed Mary,
ever a virgin, Blessed Michael the Archangel \dots\ to pray to the Lord
our God for me'' (Cummiskey 1861, pp.~xvi--xvii). The confession places
the sinner before the one who cast down the first rebel and asks the
victor over pride to pray for the man who has sinned by pride. Michael
stands in the list as the prince of the heavenly host, as the Office
calls him, \latin{princeps militiae Angelorum} (\emph{BR}, 29 Sept.,
fourth responsory).

\paragraph{The incense and the angel of the altar.} At a solemn Mass the
priest blesses the incense in Michael's name:
\begin{quote}
\latin{Per intercessionem beati Michaelis archangeli stantis a dexteris
Altaris Incensi, et omnium electorum suorum, incensum istud dignetur
Dominus benedicere, et in odorem suavitatis accipere. Per Christum Dominum
nostrum. Amen.} (Cummiskey 1861, p.~xxiv)
\end{quote}
``May the Lord, by the intercession of blessed Michael the archangel,
standing at the right hand of the Altar of Incense, and of all his elect,
vouchsafe to bless this incense, and receive it as an odour of
sweetness'' (ibid.). Two scenes of Scripture stand behind the prayer. In
the Temple, while Zachary offered incense, ``there appeared to him an
angel of the Lord, standing on the right side of the altar of incense''
(Luke 1:11). In the heavenly liturgy of the Apocalypse ``another angel
came and stood before the altar, having a golden censer: and there was
given to him much incense, that he should offer of the prayers of all
saints, upon the golden altar which is before the throne of God. And the
smoke of the incense of the prayers of the saints ascended up before God
from the hand of the angel'' (Apoc 8:3--4). The Church reads her own
incense as that prayer and asks the Archangel who stands at the altar to
carry it. The same verses furnish the Offertory of the Masses of Saint
Michael, Saint Gabriel, and Saint Raphael in the pre-1955 and 1962
Missals, \latin{Stetit Angelus juxta aram templi, habens thuribulum
aureum in manu sua} (\emph{MR}~1920, Mame 1922, p.~592), and recur
through the Office of the angels as versicle and responsory (\emph{BR},
29 Sept.\ and 2 Oct.). The rite of incense teaches in act what Hilary
teaches in the lesson of the Guardian Angels: the angels offer to God the
prayers of the saved (below, on 2 October).

\paragraph{The Preface and the Sanctus.} At the Preface the Church says
aloud that she prays with the whole court of heaven. The Common Preface
names the orders as it goes:
\begin{quote}
\latin{Per quem majestatem tuam laudant angeli, adorant dominationes,
tremunt potestates; Coeli coelorumque virtutes, ac beata seraphim, socia
exultatione concelebrant. Cum quibus et nostras voces, ut admitti jubeas
deprecamur, supplici confessione dicentes:} (Cummiskey 1861, p.~xxviii)
\end{quote}
``Thro' Christ our Lord: by whom the angels praise thy majesty, the
dominations adore it, the powers tremble before it, the heavens, the
heavenly virtues, and blessed seraphim, with common jubilee glorify it.
Together with whom we beseech thee, that we may be admitted to join our
humble voices'' (ibid.). The angels praise the Father ``through Christ'':
the Preface puts the whole heavenly worship under the mediation of the
Son, which is the doctrine of the Catechism that Christ is the centre of
the angelic world (CCC~331; Section~\ref{sec:faith}). The Preface of the
Trinity, said on the Sundays that have no proper Preface, ends
\latin{Quam laudant angeli atque archangeli, cherubim quoque ac seraphim:
qui non cessant clamare quotidie, una voce dicentes}, ``Whom the angels
and archangels, the cherubim also and seraphim praise: and cease not
daily to cry out with one voice saying'' (Cummiskey 1861, pp.~xxviii--xxix).
The Preface of the Nativity sings \latin{cum angelis et archangelis, cum
thronis et dominationibus; cumque omni militia coelestis exercitus},
``with the angels and archangels, with the thrones and dominations; and
with all the heavenly host'' (ibid., pp.~xxix--xxx), and the Preface of
Pentecost, \latin{supernae virtutes atque angelicae potestates}, ``the
heavenly virtues also and all the angelic powers'' (ibid.,
pp.~xxxi--xxxii). Between them these Prefaces name the Angels,
Archangels, Thrones, Dominations, Virtues, Powers, Cherubim, and
Seraphim: the Church's thanksgiving calls the hierarchy of the Fathers
to prayer order by order.

Every Preface ends in the Sanctus, the song Isaias heard: ``Holy, holy,
holy, the Lord God of hosts, all the earth is full of his glory'' (Isa
6:3), and John heard again before the throne: ``Holy, Holy, Holy, Lord God
Almighty, who was and who is and who is to come'' (Apoc 4:8). The Church
sings \latin{Sanctus, sanctus, sanctus, Dominus Deus Sabaoth. Pleni sunt
coeli et terra gloria tua}, ``Holy, holy, holy, Lord God of Sabaoth.
Heaven and earth are full of thy glory'' (Cummiskey 1861, p.~xxviii).
Cyril of Jerusalem explained this moment of the liturgy to the newly
baptized. The bishop, he says, makes mention ``of Angels, Archangels,
Virtues, Dominions, Principalities, Powers, Thrones; of the Cherubim with
many faces,'' and then:
\begin{quote}
We make mention also of the Seraphim, whom Esaias in the Holy Spirit saw
standing around the throne of God, and with two of their wings veiling
their face, and with twain their feet, while with twain they did fly,
crying Holy, Holy, Holy, is the Lord of Sabaoth. For the reason of our
reciting this confession of God, delivered down to us from the Seraphim,
is this, that so we may be partakers with the hosts of the world above in
their Hymn of praise. (Cyril, \emph{Catechetical Lectures} 23
[\emph{Mystagogical} 5].6)
\end{quote}
Aquinas gives the same reading in the order of the Mass: after the
Preface ``the people devoutly praise Christ's Godhead, saying with the
angels: `Holy, Holy, Holy'; and His humanity, saying with the children:
`Blessed is he that cometh''\,'' (\ST{III q.~83 a.~4 co.}). The Sanctus
is the angels' confession of the Godhead, received from the Seraphim and
given back to God by men who are to share the Seraphim's praise. The
postconciliar Order of Mass keeps it: in the Directory of 2001 the
Congregation for Divine Worship teaches that ``in the celebration of the
sacred mysteries, the Church associates herself with the angelic hymn and
proclaims the thrice holy God (cf.\ Isaiah 6,\,3)'' (\emph{Directory on
Popular Piety} 215).

\paragraph{The angel of the sacrifice.} After the consecration the Roman
Canon asks that the offering be carried by an angel to the altar on high:
\begin{quote}
\latin{Supplices te rogamus, omnipotens Deus: jube haec perferri per
manus sancti Angeli tui in sublime altare tuum, in conspectu divinae
majestatis tuae: ut, quotquot, ex hac altaris participatione sacrosanctum
Filii tui Corpus, et Sanguinem sumpserimus, omni benedictione caelesti et
gratia repleamur.} (\emph{MR}~1920, Canon Missae)
\end{quote}
The Latin says ``thy holy Angel,'' in the singular; Cummiskey's English
reads ``command these things to be carried by the hands of thy holy
angels to thy altar on high, in the sight of thy divine Majesty''
(Cummiskey 1861, pp.~xxxvii--xxxviii). Aquinas meets the objection that
the prayer asks a change of place for Christ's body, which neither comes
into the sacrament nor leaves it by local motion. The priest does not ask
that the sacramental species or Christ's true body be borne up,
\begin{quote}
but he offers this prayer for Christ's mystical body, which is signified
in this sacrament, that the angel standing by at the Divine mysteries may
present to God the prayers of both priest and people, according to
Apocalypse 8:4: ``And the smoke of the incense of the prayers of the
saints ascended up before God, from the hand of the angel.'' But God's
``altar on high'' means either the Church triumphant, unto which we pray
to be translated, or else God Himself, in Whom we ask to share \dots\ Or
else by the angel we are to understand Christ Himself, Who is the ``Angel
of great counsel'' (Isaiah 9:6: Septuagint), Who unites His mystical body
with God the Father and the Church triumphant. (\ST{III q.~83 a.~4 ad~9})
\end{quote}
Aquinas adds that the Mass takes its name from this sending: \latin{missa},
``because the priest sends [\latin{mittit}] his prayers up to God through
the angel, as the people do through the priest'' (ibid.). Both readings
stand in the \emph{Summa}: an angel assisting at the altar carries the
Church's prayer, or the Angel of great counsel carries his own Body into
the presence of the Father. The Directory of 2001 cites the Supplices
among the prayers of the postconciliar Missal in which the Church invokes
the angels' ``assistance so that the Eucharistic sacrifice `may be taken
[to your] altar in heaven, in the presence of [\dots] divine majesty'\,''
(\emph{Directory} 215).

Gregory the Great had taught the same presence of the angels in the hour
of sacrifice. In the fourth book of the \emph{Dialogues} he writes:
\begin{quote}
\dots\ for what right believing Christian can doubt, that in the very hour of
the sacrifice, at the words of the Priest, the heavens be opened, and the
quires of Angels are present in that mystery of Jesus Christ; that high
things are accompanied with low, and earthly joined to heavenly, and that
one thing is made of visible and invisible? (Gregory, \emph{Dialogues}
IV.58, English of 1608, ed.\ Gardner, 1911, p.~256)
\end{quote}
The Mass is heaven opened. The angels who surround the throne in the
Apocalypse surround the altar at the words of consecration, and the
Church, which cannot see them, confesses their presence in her Preface
and her Canon.

\paragraph{The Litany of the Saints.} When the Church calls on the whole
Church triumphant, at the blessing of the font on Holy Saturday and in
her rites of supplication, the angels follow the Mother of God:
\latin{Sancte Michael, ora. Sancte Gabriel, ora. Sancte Raphael, ora.
Omnes sancti Angeli et Archangeli, orate pro nobis. Omnes sancti beatorum
Spirituum ordines, orate} (\emph{MR}~1920, Sabbato Sancto, the litany).
The three angels whom Scripture names are invoked by name; all the others
are invoked by order. The litany keeps the rule by which the Church names
the angels (Section~\ref{sec:devotion}; Appendix~\ref{app:names}).

\sectionguard
\subsection{The Paschal Vigil and the Office of the Night}

\paragraph{The Exsultet.} The deacon's praise of the paschal candle opens
by summoning the angels. In the pre-1955 Holy Saturday he sings,
\latin{Exsultet jam Angelica turba coelorum: exsultent divina mysteria:
et pro tanti Regis victoria tuba insonet salutaris} (\emph{MR}~1920,
Sabbato Sancto, \emph{Benedictio Cerei}, p.~233), ``Let now the heavenly
troop of angels rejoice'' (Cummiskey, \emph{Roman Missal}, Philadelphia
1843, Holy Saturday). The first to be called to the joy of the
resurrection are the angels, who announced the Rising at the tomb and
whose voice the Apostle hears at the Lord's return: ``the Lord himself
shall come down from heaven with commandment and with the voice of an
archangel and with the trumpet of God'' (1 Thess 4:15). The Exsultet
unites heaven, earth, and Mother Church in one proclamation, and heaven's
part is sung first.

\paragraph{The angels in the house at night.} At the end of every day,
in the collect of Compline, the Church asks that the angels keep her
houses:
\begin{quote}
\latin{Visita, quaesumus Domine, habitationem istam, et omnes insidias
inimici ab ea longe repelle: Angeli tui sancti habitent in ea, qui nos in
pace custodiant; et benedictio tua sit super nos semper. Per Dominum.}
(\emph{BR}, \emph{Pars autumnalis}, Ad Completorium, p.~151)
\end{quote}
``Visit, we beseech Thee, O Lord, this habitation, and drive far from it
all snares of the enemy: let Thine holy Angels dwell herein, to keep us in
peace, and may Thy blessing be always upon us'' (Bute, IV, Compline). The
prayer sets the two ministries of the angels side by side as the treatise
does: the snares of the enemy are driven off and the holy angels dwell
with men to guard them (\ST{I qq.~113--114}; Sections~\ref{sec:guardians},
\ref{sec:assaults}). The Directory reads the prayer in the same way: the
Church ``implores God to send his Angels at the end of the day to protect
the faithful as they sleep'' (\emph{Directory} 215).

\paragraph{``In the sight of the angels.''} The versicle that closes the
hymns of every Office of the angels is the psalmist's, \latin{In conspectu
Angelorum psallam tibi Deus meus. Adorabo ad templum sanctum tuum, et
confitebor nomini tuo} (\emph{BR}, 29 Sept.\ and 2 Oct.): ``I will sing
praise to thee in the sight of the angels'' (Ps 137[138]:1). The Directory
takes the verse as the rule of all the Church's morning prayer: ``The
office of lauds is celebrated in their presence (cf.\ Ps 137,\,1)''
(\emph{Directory} 215). The Church does not pray alone at any hour; she
prays before an audience of angels.

\sectionguard
\subsection{The Dying and the Dead}

The liturgy gives the angels their largest office at death. Scripture
shows the pattern once and for all: ``the beggar died and was carried by
the angels into Abraham's bosom'' (Luke 16:22). The Roman rites of the
dying, of the Requiem, and of burial pray that the pattern be fulfilled in
every Christian death.

\paragraph{The commendation of the soul.} When the dying man is in his
agony the priest sends him forth in the name of the Trinity and of the
whole heavenly court:
\begin{quote}
\latin{Proficiscere, anima christiana, de hoc mundo, in nomine Dei Patris
omnipotentis, qui te creavit: in nomine Jesu Christi Filii Dei vivi, qui
pro te passus est: in nomine Spiritus Sancti, qui in te effusus est: in
nomine Angelorum et Archangelorum: in nomine Thronorum et Dominationum: in
nomine Principatuum et Potestatum: in nomine Cherubim et Seraphim: in
nomine Patriarcharum et Prophetarum \dots} (\emph{RR}~1872,
\emph{Ordo commendationis animae}, pp.~115--116)
\end{quote}
The soul is sent out of the world in the name of the three Persons, and
then in the name of the angelic orders, pair by pair, before the
patriarchs, apostles, martyrs, and virgins. The next prayer asks that as
the soul goes out of the body ``the shining company of angels'' come to
meet it, \latin{splendidus Angelorum coetus occurrat} (ibid., p.~116).
The man who dies in the Church dies into the company of the angels.

\paragraph{The Subvenite.} As soon as the soul has left the body the
Church says the responsory that asks the saints and angels to receive it:
\begin{quote}
\latin{Subvenite Sancti Dei, occurrite Angeli Domini, Suscipientes animam
ejus, Offerentes eam in conspectu Altissimi. V.~Suscipiat te Christus, qui
vocavit te, et in sinum Abrahae Angeli deducant te.} (\emph{RR}~1872,
pp.~131--132)
\end{quote}
``Come to his assistance, ye Saints of God, come forth to meet him, ye
Angels of the Lord: Receiving his soul: Offering it in the sight of the
Most High. V.~May Christ receive thee, who hath called thee, and may the
Angels bear thee into Abraham's bosom'' (\emph{Manual of Prayers}, 1889,
pp.~531--532). The angels do for the Christian what they did for Lazarus;
and the verb the Church gives them, \latin{offerentes}, is the verb of
sacrifice: they present the soul before God as the angel of the
Supplices presents the Church's offering.

\paragraph{Michael the standard-bearer.} The Offertory of the Mass for
the Dead names the angel who leads the souls:
\begin{quote}
\latin{Domine Jesu Christe, Rex gloriae, libera animas omnium fidelium
defunctorum de poenis inferni et de profundo lacu: libera eas de ore
leonis, ne absorbeat eas tartarus, ne cadant in obscurum: sed signifer
sanctus Michael repraesentet eas in lucem sanctam: Quam olim Abrahae
promisisti et semini ejus.} (\emph{MR}~1920, \emph{Missae pro
defunctis})
\end{quote}
``\dots\ and let the standard bearer, St Michael, bring them into the holy
light: Which thou promisedst of old to Abraham and his posterity''
(Cummiskey 1843, Masses for the Dead). The Office of Michael sings the
same ministry. At Lauds the Church gives God's own word to the
Archangel, \latin{Archangele Michael, constitui te principem super omnes
animas suscipiendas}, ``Michael Mine Archangel, I have appointed thee for
a prince over the ingathering of souls''; at Matins she sings
\latin{Venit Michael Archangelus cum multitudine Angelorum, cui tradidit
Deus animas sanctorum, Ut perducat eas in paradisum exsultationis},
rendered by Bute ``Where Angels lead the spirits of the blessed dead the
glad procession moves with Michael at his head, to lead them into the
garden of Eden''; and in the third nocturn, \latin{Angelus Archangelus
Michael, Dei nuntius pro animabus justis}, ``he is the messenger whom God
sendeth to all the souls of the righteous'' (\emph{BR}, 29 Sept.,
pp.~381--384; Bute, IV, pp.~595--598). The \latin{signifer} of the
Requiem is the \latin{salutis signifer} of the Vespers hymn of his feast,
who as victor unfurls the cross, \latin{explicat victor crucem} (below, 29
September): the standard he carries before the souls is the cross of
Christ.

\paragraph{In paradisum.} As the body is carried from the church to the
grave the Church sings:
\begin{quote}
\latin{In paradisum deducant te Angeli: in tuo adventu suscipiant te
Martyres, et perducant te in civitatem sanctam Jerusalem. Chorus Angelorum
te suscipiat, et cum Lazaro quondam paupere aeternam habeas requiem.}
(\emph{RR}~1872, \emph{Exsequiarum ordo}, p.~141)
\end{quote}
``May the Angels lead thee into Paradise; at thy coming may the Martyrs
receive thee, and bring thee into the holy City, Jerusalem. May the Choir
of Angels receive thee, and with Lazarus, once a beggar, mayest thou have
eternal rest'' (\emph{Manual of Prayers}, 1889, pp.~583--584). The
antiphon ends where the Gospel began, with Lazarus carried by the angels.
The Directory sums up these rites: the Church ``prays that the celestial
spirits come to the assistance of the faithful in their last agony, and
in the rite of obsequies, invokes God to send his Angels to accompany the
souls of just into paradise'' (\emph{Directory} 215). The common teaching
that the angels assist the dying and bring the just to judgment and to
glory is here the Church's prayer at every Christian death.

\sectionguard
\subsection{The Feasts of the Angels}

The Roman calendar before 1955 keeps five feasts of the angels: the
Dedication of Saint Michael on 29 September, his Apparition on 8 May, the
Guardian Angels on 2 October, Saint Gabriel on 24 March, and Saint
Raphael on 24 October. The 1962 Missal keeps four in the general calendar
and prints the Apparition among the Masses for particular places. The
postconciliar calendar keeps the Feast of Saints Michael, Gabriel, and
Raphael on 29 September and the Memorial of the Guardian Angels on 2
October. Every feast of the angels in the Tridentine Missal opens with
the same Introit, from the psalm that bids the angels bless the Lord:
\latin{Benedicite Dominum, omnes Angeli ejus: potentes virtute, qui
facitis verbum ejus, ad audiendam vocem sermonum ejus}, ``Bless the Lord,
all ye his angels: you that are mighty in strength, and execute his word,
hearkening to the voice of his orders'' (Ps 102[103]:20; \emph{MR}~1920,
8 May, 2 Oct., 24 Mar., 24 Oct.). The Church begins each feast by
joining the angels in their own work, which is to bless God and to do his
word.

\subsubsection{The Dedication of Saint Michael (29 September)}

\paragraph{The day and its rank.} The feast is the oldest of the angels
in the Roman books. The Leonine Sacramentary, the oldest of the Roman
sacramentaries, keeps on the day before the Kalends of October the
\latin{natale basilicae Angeli in Salaria}, the anniversary of the church
of the Angel on the Salarian Way, and its Preface gives the reason for
such a feast:
\begin{quote}
\latin{Vere dignum \dots\ in die festivitatis hodiernae quo in honorem
beati archangeli Michaelis sacrata nomini tuo loca Divinis sunt instituta
mysteriis: quamvis enim illius sublimis gloriosaeque substantiae sit
habitatio semper in caelis tuorum tamen fidelium praesumit affectus pro
tuae reverentia potestatis per haec piae devotionis officia quoddam
retinere pignus in terris adstantium in conspectu tuo iugiter
ministrorum.} (\emph{Sacramentarium Leonianum}, ed.\ Feltoe, 1896,
p.~106)
\end{quote}
Though the Archangel's dwelling is always in heaven, the faithful, out of
reverence for God's power, hold on earth a pledge of the ministers who
stand for ever in God's sight. Another Preface of the same book confesses that the
angelic nature, withdrawn from bodily sight, is seen by the gaze of
faith:
\latin{quae etsi humano generi corporeo conspectu subtrahitur fidei tamen
videtur intuitu} (ibid., p.~107). The Gregorian Sacramentary sets the
\latin{Dedicatio basilicae sancti Angeli} on the third day before the
Kalends, 29 September, with the collect the Tridentine Missal still
prays (\emph{Gregorian Sacramentary}, ed.\ Wilson, 1915, under
III~Kal.\ Oct.). The Breviary gives the feast its scope. Relating the
dedication of Saint Michael's church at Rome on the third day before the
Kalends of October, the lessons of 8 May add: \latin{quo die etiam
omnium Angelorum memoriam Ecclesia celebrat}, ``on the which day the
Church also holdeth in remembrance All Angels'' (\emph{BR}, \emph{Pars
verna}, 8 May, lesson vi, p.~572; Bute, II, p.~870). Michaelmas is the
feast of all the angels under the name of their prince.

The Breviary of 1898 and Bute's Breviary of 1908 give the day the rank of
a double of the second class (\emph{BR}, p.~378; Bute, IV, p.~592); the
Missal of the 1920 typical edition gives it the first class, \latin{Duplex
I classis} (\emph{MR}~1920, Mame 1922, kalendarium and p.~725), and the
1962 Missal keeps it of the first class (\emph{MR}~1962, p.~670). The
postconciliar calendar keeps 29 September as the Feast of Saints Michael,
Gabriel, and Raphael, Archangels.

\paragraph{The Mass.} The Mass of the day in the pre-1955 and 1962
Missals is the Mass \latin{Benedicite}, which Mame's Missal prints in
full on 8 May and appoints again on 29 September, \latin{Missa Benedicite,
ut in die 8 Maji \dots\ praeter ritum}, with the Gradual and Alleluia proper to
the autumn feast (\emph{MR}~1920, Mame 1922, pp.~592--593, 725--726). Its collect is the
Gregorian collect:
\begin{quote}
\latin{Deus, qui miro ordine Angelorum ministeria, hominumque dispensas:
concede propitius; ut a quibus tibi ministrantibus in caelo semper
assistitur, ab his in terra vita nostra muniatur. Per Dominum.}
(\emph{BR}, 29 Sept., p.~384)
\end{quote}
``O God, Who hast ordained and constituted the services of angels and men
in a wonderful order, mercifully grant that as Thy holy angels alway do
Thee service in heaven, so, by Thy appointment, they may succour and
defend us on earth'' (Bute, IV, pp.~598--599). The prayer holds the whole
theology of the angels' ministry in one sentence. There is one order,
\latin{miro ordine}, in which God disposes the offices of angels and of
men together; the angels who assist before God in heaven are the angels
who defend men's lives on earth. Aquinas teaches the same when he sets the
angels over the government of bodies and of men and makes the lower
creation receive God's providence through the higher (\ST{I
qq.~110--113}; Sections~\ref{sec:government}, \ref{sec:hierarchies}).
Isidore distinguishes the angels who minister from the angels who
assist by the two verbs of Daniel's vision, ``thousands of thousands
ministered to him, and ten thousand times a hundred thousand stood before
him'' (Dan 7:10): \latin{alii sunt qui administrant, alii qui adsistunt}
(Isidore, \emph{Etym.}\ VII.5.19). The collect asks that the spirits who
stand before God should also minister to us.

The Lesson is the opening of the Apocalypse, in which ``God \dots\
signified, sending by his angel to his servant John'' (Apoc 1:1); the
Gradual repeats Psalm 102[103]:20; the Alleluia is the Church's own
invocation, \latin{Sancte Michael Archangele, defende nos in praelio: ut
non pereamus in tremendo judicio}, ``Holy Michael the Archangel, defend
us in the battle: that we may not perish in the dreadful judgment''
(\emph{MR}~1920, Mame 1922, p.~725; Cummiskey 1843, p.~667). The battle is
the last one, the judgment at which the dragon is finally cast out; the
same petition opens the prayer to Saint Michael that Leo~XIII set after
Low Mass (Section~\ref{sec:devotion}). In paschal time a
second verse sings the terror of the Archangel's coming, \latin{Concussum
est mare et contremuit terra, ubi Archangelus Michael descendit de
caelo}, ``The sea shook, and the earth trembled, when Michael the
Archangel came down from heaven'' (\emph{MR}~1920, Mame 1922, p.~592;
Cummiskey 1843, p.~667). The Gospel is Matthew 18:1--10, which ends ``See
that you despise not one of these little ones: for I say to you, that
their angels in heaven always see the face of my Father who is in heaven''
(Matt 18:10): on the feast of the prince of the angels the Church reads
the charter of the guardian angels. The Offertory is \latin{Stetit
Angelus} (Apoc 8:3--4); the Communion is the canticle of the Three
Children, \latin{Benedicite, omnes Angeli Domini, Dominum: hymnum dicite,
et superexaltate eum in saecula} (Dan 3:58).

The Secret and the Postcommunion reach back to the Leonine Sacramentary. The Secret,
\latin{Hostias tibi, Domine, laudis offerimus, suppliciter deprecantes:
ut easdem, angelico pro nobis interveniente suffragio, et placatus
accipias, et ad salutem nostram provenire concedas} (\emph{MR}~1920,
Mame 1922, p.~592), already stands in the Leonine Sacramentary for the
church on the Salarian Way (Feltoe, p.~106) and in the Gregorian for 29
September (Wilson, p.~106). The Church offers the sacrifice of praise and
asks that it be received through the angels' pleading, \latin{angelico
pro nobis interveniente suffragio}. The Postcommunion asks that, upheld by
Michael's intercession, we may reach with the mind what we have followed
with the mouth: \latin{Beati Archangeli tui Michaelis intercessione
suffulti: supplices te, Domine, deprecamur; ut, quod ore prosequimur,
contingamus et mente} (\emph{MR}~1920, Mame 1922, p.~593). The Leonine
form asks that we may reach in mind those whom we follow with honour,
\latin{ut quos honore prosequimur contingamus et mente} (Feltoe, p.~108),
and the Gregorian keeps that reading (Wilson, p.~106). The Tridentine
prayer turns the honour paid to the angels into the Church's praise: what
we have sung with the lips we ask to touch with the mind.

\paragraph{The Office.} The Office of Michaelmas is the Church's fullest
teaching on the angels in her own voice. The Invitatory adores
\latin{Regem Archangelorum Dominum}, ``The Lord, He is the King of the
Archangels'' (\emph{BR}, p.~378; Bute, IV, p.~593). The Vespers hymn
sings Michael as the standard-bearer of salvation who unfurls the cross:
\begin{quote}
\latin{Tibi mille densa millium\\
Ducum corona militat:\\
Sed explicat victor crucem\\
Michael salutis signifer.\\
Draconis hic dirum caput\\
In ima pellit tartara,\\
Ducemque cum rebellibus\\
Caelesti ab arce fulminat.} (\emph{BR}, 29 Sept., p.~378)
\end{quote}
The English that the \emph{Raccolta} prints beside the hymn renders the
stanzas: ``Thy thousand, thousand hosts are spread, / Embattled o'er the
azure sky; / But Michael bears thy standard dread, / And lifts the mighty
Cross on high. / He in that Sign the rebel powers /
Did with their Dragon Prince expel; / And hurled them from the Heaven's
high towers, / Down like a thunderbolt to hell'' (\emph{Raccolta}, 1910,
n.~289, p.~259). The hymn gives the doctrine of the fall in one image:
the victory over the dragon is won under the sign of the cross. The
Lauds hymn \latin{Christe, sanctorum decus Angelorum} gives each of the
three named archangels his office:
\begin{quote}
\latin{Angelus pacis Michael in aedes\\
Caelitus nostras veniat, serenae\\
Auctor ut pacis lacrymosa in orcum\\
Bella releget.\\
Angelus fortis Gabriel, ut hostes\\
Pellat antiquos, et amica caelo,\\
Quae triumphator statuit per orbem,\\
Templa revisat.\\
Angelus nostrae medicus salutis,\\
Adsit e caelo Raphael, ut omnes\\
Sanet aegrotos, dubiosque vitae\\
Dirigat actus.} (\emph{BR}, 29 Sept., p.~384)
\end{quote}
Copeland's English in Bute's Breviary: ``Angel of peace, may Michael to
our dwelling / Down from high Heaven in mighty calmness come, / Breathing
serenest peace, wild war dispelling / With all her sorrows to the infernal
gloom. / Angel of might, may Gabriel swift descending, / Far from our
gates our ancient foes repel, / And his own triumphs o'er the world
defending, / In temples dear to Heaven return and dwell. / Angel of
health, may Raphael lighten o'er us, / To every sick bed speed his healing
flight, / In times of doubt direct the way before us, / And through life's
mazes guide our steps aright'' (Bute, IV, p.~598). The hymn gives Michael peace,
Gabriel strength, and Raphael healing; the last two answer to the names as
Gregory interprets them in the lesson of the day, ``the Strength of
God'' and ``the Medicine-of-God'' (Bute, IV, p.~596;
Appendix~\ref{app:names}).

The lessons of the first nocturn are Daniel's two visions of the court of
heaven and of Michael, ``one of the chief princes'' (Dan 7:9--10;
10:4--14; \emph{BR}, pp.~379--380). The second nocturn reads Gregory the
Great's thirty-fourth homily on the Gospels, the sermon in which he
gathered the nine orders from Scripture. The fourth lesson gives the
count:
\begin{quote}
\latin{Novem Angelorum ordines dicimus, quia videlicet esse, testante
sacro eloquio, scimus: Angelos, Archangelos, Virtutes, Potestates,
Principatus, Dominationes, Thronos, Cherubim atque Seraphim.}
(\emph{BR}, pp.~380--381; Gregory, \emph{Hom.\ in Ev.} 34.7)
\end{quote}
``We say that there are nine Orders of Angels, for, by the witness of the
holy Word, we know that there be Angels, Archangels, Mights, Powers,
Principalities, Dominions, Thrones, Cherubim, and Seraphim'' (Bute, IV,
pp.~594--595). Gregory adds four from Ephesians, a fifth, the Thrones,
from Colossians, and Angels, Archangels, Cherubim, and Seraphim from the
rest of Scripture, ``and we find that the Orders of Angels are beyond all
doubt nine'' (ibid.). The fifth lesson gives Gregory's teaching on the
name: ``the word `Angel' is the designation, not of a nature, but of an
office,'' \latin{Angelorum vocabulum nomen est officii, non naturae}; the
holy spirits ``are Angels only when they are sent as Messengers,'' and
``they who go on the lesser messages are called Angels: they who go on
the greater Archangels. Hence it is that unto the Virgin Mary was sent no
common Angel, but the Archangel Gabriel'' (Bute, IV, p.~595; \emph{BR},
p.~381; Gregory, \emph{Hom.\ in Ev.}\ 34.8). The sixth lesson interprets the
names: whenever a work of wondrous power is to be done Michael is sent,
``that by that very thing, and by his name, we may remember that none is
able to do as God doeth'' (Bute, IV, p.~596; Gregory, \emph{Hom.\ in Ev.}\ 34.9).
Thus on the feast of the angels the Church reads to her clergy the
patristic source of the doctrine of the nine choirs that Aquinas received
beside Dionysius (\ST{I q.~108 a.~6}; Section~\ref{sec:hierarchies};
Appendix~\ref{app:orders}).

The third nocturn reads Jerome on Matthew 18. His ninth lesson is the
classic patristic sentence on the guardian angels:
\begin{quote}
\latin{Magna dignitas animarum, ut unaquaeque habeat ab ortu nativitatis
in custodiam sui Angelum delegatum.} (\emph{BR}, 29 Sept., p.~383;
Jerome, \emph{In Matth.} III, on 18:10)
\end{quote}
``Oh, how great is the dignity of souls, whereof every one hath from its
birth an Angel appointed to guard it!'' (Bute, IV, p.~597). Jerome finds
the same ministry in the angels of the seven churches of the Apocalypse
and in the Apostle's word that women veil their heads in church ``because
of the Angels'' (ibid.). Aquinas cites the sentence for the question
whether each man has a guardian from his birth (\ST{I q.~113 a.~5};
Section~\ref{sec:guardians}).

At Lauds the antiphons pass from the altar of incense to the war in
heaven, to Michael's charge over souls, and to the whole hierarchy
bidden to praise: \latin{Angeli, Archangeli, Throni et Dominationes,
Principatus et Potestates, Virtutes caelorum, laudate Dominum de caelis},
``O ye Angels and Archangels, O ye Thrones and Dominions, O ye
Principalities and Powers, O ye mighty Ones of heaven, praise ye the Lord
from the heavens!'' (\emph{BR}, p.~384; Bute, IV, p.~598). The Office
closes at Second Vespers with the antiphon that the faithful also pray in
private (Section~\ref{sec:devotion}):
\begin{quote}
\latin{Princeps gloriosissime, Michael Archangele, esto memor nostri: hic
et ubique semper precare pro nobis Filium Dei, alleluia, alleluia.}
(\emph{BR}, p.~385)
\end{quote}
``O thou Prince most glorious, Michael the Archangel, remember us---and
here, and everywhere, alway entreat for us the countenance of the Son of
God'' (Bute, IV, p.~600). The Archangel is asked to pray; the one he is
asked to pray to is the Son.

The Martyrology read at Prime on the eve gives the reason of the day: ``On the morrow is made the worshipful Commemoration of
the Blessed Archangel Michael; on the which day was consecrated in his
name the church upon Mount Gargano, mean as to building, but filled with
power from heaven'' (Bute, IV, pp.~591--592). The shrine of Gargano and the
other sanctuaries of the Archangel are related in
Section~\ref{sec:devotion}.

\paragraph{The postconciliar feast.} Since the reform of the calendar the
Roman Rite keeps on 29 September the Feast of Saints Michael, Gabriel,
and Raphael, Archangels, and gathers into it the separate feasts of
Gabriel and Raphael. The \emph{Ordo lectionum Missae} of 1981 appoints for
the first reading either Daniel's vision of the Ancient of days, before
whom ``thousands of thousands ministered'' (Dan 7:9--10, 13--14), or the
war in heaven in which ``Michael and his angels fought with the dragon''
(Apoc 12:7--12a); for the psalm, the verse that the Tridentine Office sang
as its versicle, ``I will sing praise to thee in the sight of the
angels'' (Ps 137[138]:1--5); for the verse before the Gospel, ``Bless the
Lord, all ye his hosts'' (Ps 102[103]:21); and for the Gospel, the promise
to Nathanael, ``you shall see the heaven opened and the angels of God
ascending and descending upon the Son of man'' (John 1:47--51)
(\emph{Ordo lectionum Missae} 1981, n.~647). The Gospel sets the whole
feast in Christ: the angels go up and come down upon the Son of man, the
ladder of Jacob's dream, and their ministry to men is his. The
Catechism teaches the same: ``Christ is the centre of the
angelic world. They are his angels'' (CCC~331).

\subsubsection{The Apparition of Saint Michael (8 May)}

The pre-1955 Missal and Breviary keep on 8 May a second feast of the
Archangel, \latin{In Apparitione S.~Michaelis Archangeli}, a greater
double (\emph{MR}~1920, Vatican printing, kalendarium and proper of 8 May;
Bute, II, p.~867). Its Mass
is the Mass \latin{Benedicite} with the alleluias of Eastertide
(\emph{MR}~1920, Mame 1922, pp.~592--593), and Mame's Missal directs that
a votive Mass of Saint Michael take this form in paschal time and the form
of 29 September outside it (ibid., p.~593). The 1962 Missal prints the
day among the Masses for particular places: \latin{Eodem die 8 maii. In
Apparitione S.~Michaelis Archangeli. Missa Benedicite, ut in Missali, die
29 septembris} (\emph{MR}~1962, \emph{Missae pro aliquibus locis},
p.~[161]). The Office of the day reads in its second nocturn the relation
of the Archangel's appearing on Mount Gargano in the pontificate of
Gelasius~I; the lessons are given with the sanctuaries of the Archangel in
Section~\ref{sec:devotion}.

\subsubsection{The Holy Guardian Angels (2 October)}

\paragraph{The day and its rank.} The Guardian Angels were honoured in
many churches before the feast was made universal. \emph{The Liturgical
Year} of Guéranger's continuators relates that Paul~V permitted the feast
in 1608 and that Clement~X made it of precept on 2 October, the first
free day after Michaelmas, in 1670 (\emph{The Liturgical Year}, Time after
Pentecost V, English of the Stanbrook Benedictines, October 2). The
Breviary of 1898 and the Missal of the 1920 typical edition keep it as a
greater double (\emph{BR}, p.~397; \emph{MR}~1920, Mame 1922, kalendarium
and p.~726); the 1962 Missal as a feast of the
third class (\emph{MR}~1962, p.~672); the postconciliar calendar as a
Memorial.

\paragraph{The Mass.} The collect states the doctrine of the day:
\begin{quote}
\latin{Deus, qui ineffabili providentia sanctos Angelos tuos ad nostram
custodiam mittere dignaris: largire supplicibus tuis; et eorum semper
protectione defendi, et aeterna societate gaudere. Per Dominum.}
(\emph{BR}, 2 Oct., p.~403)
\end{quote}
``O God, Who in Thine unspeakable Providence hast been pleased to give
Thine holy Angels charge over us, to keep us, mercifully grant unto our
prayers, that we be both ever fenced by their wardship here, and
everlastingly blessed by their fellowship hereafter'' (Bute, IV, p.~619).
The prayer names the two gifts of the guardian angel: protection now and
companionship for ever. The guardians are sent by an ``unspeakable
providence,'' and the men they guard are called to be their fellow
citizens.

The Lesson is God's promise in Exodus, ``Behold I will send my angel, who
shall go before thee, and keep thee in thy journey, and bring thee into
the place that I have prepared'' (Exod 23:20--23). The Gradual is the
psalm the tempter quoted and the Church sings in its true sense: ``For he
hath given his angels charge over thee; to keep thee in all thy ways. In
their hands they shall bear thee up: lest thou dash thy foot against a
stone'' (Ps 90[91]:11--12). The Gospel is again Matthew 18:1--10, whose
last verse gives the feast its Scripture; the Offertory is the psalm of
the Introit (Ps 102[103]:20--21); the Communion is the canticle of the
Three Children (Dan 3:58) (\emph{MR}~1920, Mame 1922, pp.~726--729). The
Secret asks \latin{ut, perpetuis eorum praesidiis, a praesentibus periculis
liberemur, et ad vitam perveniamus aeternam}, that by their unfailing
protection we may be delivered from present dangers and come to eternal
life; the Postcommunion, rejoicing in their feast, asks \latin{ut eorum
protectione ab hostium jugiter liberemur insidiis, et contra omnia adversa
muniamur}, that by their protection we may be delivered from the snares
of our enemies and fortified against every adversity (ibid., p.~729).

\paragraph{The Office.} The Vespers hymn \latin{Custodes hominum psallimus
Angelos} sings the guardians as the companions given to fallen nature
against the envy of the fallen angel:
\begin{quote}
\latin{Custodes hominum psallimus Angelos,\\
Naturae fragili quos Pater addidit\\
Caelestis comites, insidiantibus\\
Ne succumberet hostibus.} (\emph{BR}, 2 Oct., p.~397)
\end{quote}
``Praise we those ministers celestial / Whom the dread Father chose / To be
defenders of our nature frail, / Against our scheming foes. / For, since
that from his glory in the skies / Th' Apostate Angel fell, / Burning with
envy, evermore he tries / To drown our souls in Hell'' (Caswall's
English, Bute, IV, p.~613). The Lauds hymn asks that the angel chosen
to guard us be present, \latin{Tuusque nobis Angelus / Electus ad
custodiam, / Hic adsit}; in Caswall's English, ``May Thy Guardian Angel
nigh / Keep us from all sin this day. / May he crush the deadly wiles /
Of the envious Serpent's art,'' and drive ``Plague and famine'' far
(\emph{BR}, p.~402; Bute, IV, p.~619).

The second nocturn reads Bernard's sermon on the ninetieth psalm. He
begins with the verse the Gradual sings and asks four questions of it:
``For who hath given charge? And what charge? Unto whom? And over whom?''
(Bute, IV, p.~616). The answer is the wonder of the day: ``The Highest
Majesty, therefore, hath given charge unto Angels, even His Angels. Unto
these beings so excellently exalted, so blessed, so near to Himself, even
as His own household, unto these hath He given charge over thee. Who art
thou?'' (ibid.). The fifth lesson draws the conclusion, \latin{Reverentiam
pro praesentia, devotionem pro benevolentia, fiduciam pro custodia}
(\emph{BR}, 2 Oct., p.~400; the sermon is expounded at
\S\ref{sec:fa-bernard}): ``What respect, what thankfulness, what trust, ought this word to work in
thee! Respect for their presence, thankfulness for their kindness, trust
in their safe-keeping. Walk carefully, as one with whom are Angels, as
hath been laid in charge upon them, in all thy ways. In every lodging, in
every nook, have reverence for thine Angel. Dare not to do in his
presence what thou wouldst not dare to do in mine'' (Bute, IV, p.~616).
The sixth lesson turns reverence into love and confidence: ``Let us also,
brethren, dearly love His Angels, as them with whom we are one day to be
co-heirs, and who in the meanwhile are leaders and guardians set over us
by the Father. With such guardians, whereof shall we be afraid? They that
keep us in all our ways, can neither be conquered nor corrupted, far less
can they corrupt. They are trusty, they are wary, they are mighty''
(Bute, IV, pp.~616--617; \emph{BR}, p.~400).

The third nocturn reads Hilary on Matthew 18. His ninth lesson makes the
guardians the bearers of the prayers of the faithful:
\begin{quote}
\latin{Pusillorum enim Angeli quotidie Deum vident: quia Filius hominis
venit salvare quae perdita sunt. Ergo et Filius hominis salvat, et Deum
Angeli vident, et Angeli pusillorum praesunt fidelium orationibus.
Praeesse Angelos absoluta auctoritas est. Salvatorum igitur per Christum
orationes Angeli quotidie Deo offerunt.} (\emph{BR}, 2 Oct., p.~402)
\end{quote}
``From these words we see, first, that the Son of Man saveth; secondly,
that the Angels do see God; and thirdly, that the Angels of these little
ones have the wardship over the prayers of the faithful. That the Angels
have this wardship is taught us absolutely. The Angels therefore do every
day offer to God the prayers, which they which are saved, do make to Him
in the Name of Christ'' (Bute, IV, p.~618). Hilary's lesson is the
patristic commentary on the incense of the Offertory and the Supplices of
the Canon: the angels offer daily what the Church offers at every Mass.

\paragraph{The postconciliar Memorial.} The \emph{Ordo lectionum Missae}
keeps the Tridentine readings: Exodus 23:20--23, Psalm 90[91], the verse
``Bless the Lord, all ye his hosts'' (Ps 102[103]:21), and Matthew
18:1--5, 10 (\emph{Ordo lectionum Missae} 1981, n.~650). The doctrine of
the day stands in the Catechism: ``Beside each believer stands an angel as
protector and shepherd leading him to life'' (CCC~336, citing Basil).

\subsubsection{Saint Gabriel (24 March)}

The Missal of the 1920 typical edition keeps the feast of Saint Gabriel
the Archangel on 24 March, the eve of the Annunciation, as a greater
double (\emph{MR}~1920, Mame 1922, kalendarium and pp.~552--554); the 1962
Missal as a feast of the third class (\emph{MR}~1962, p.~493); the
postconciliar calendar honours Gabriel on 29 September. The collect gives
the reason of the day:
\begin{quote}
\latin{Deus, qui inter ceteros Angelos, ad annuntiandum incarnationis tuae
mysterium, Gabrielem Archangelum elegisti: concede propitius; ut, qui
festum ejus celebramus in terris, ipsius patrocinium sentiamus in caelis.}
(\emph{MR}~1920, Mame 1922, p.~552)
\end{quote}
God chose Gabriel among all the angels to announce the mystery of the
Incarnation, and the Church asks that those who keep his feast on earth
may know his patronage in heaven. The Lesson is Daniel's second meeting
with Gabriel, ``flying swiftly'' at the time of the evening sacrifice to
reveal the weeks ``unto Christ, the prince'' (Dan 9:21--26); the
Gospel is the Annunciation (Luke 1:26--38); in Lent the Tract sings the
angel's salutation, \latin{Ave, Maria, gratia plena: Dominus tecum}
(Luke 1:28, 42, 31, 35). The Secret asks that the prayer of the blessed
Archangel Gabriel make the offering acceptable, \latin{ut, qui a nobis
veneratur in terris, sit apud te pro nobis advocatus in caelis}, that he
who is venerated by us on earth may be our advocate with God in heaven;
the Postcommunion prays \latin{ut, sicut Gabriele nuntiante,
incarnationem tuam cognovimus; ita, ipso adjuvante, incarnationis ejusdem
beneficia consequamur}, that as by Gabriel's message we have known the
Incarnation, so by his help we may obtain its benefits (ibid.,
pp.~553--554). Gregory's homily, read on 29 September, gives the reason
the Archangel of the Annunciation is ``the Strength of God'': ``he came
to herald the appearing of Him Who was content to appear lowly that He
might fight down the powers of the air'' (Bute, IV, p.~596; Gregory,
\emph{Hom.\ in Ev.}\ 34.9). The feast of Gabriel is the vigil of the mystery he
announced.

\subsubsection{Saint Raphael (24 October)}

The Missal of the 1920 typical edition keeps Saint Raphael the Archangel
on 24 October as a greater double (\emph{MR}~1920, Mame 1922, kalendarium
and pp.~741--743); Bute's Breviary for England already kept the day with
that rank in 1908 (Bute, IV, kalendar and pp.~684--686); the 1962 Missal
keeps it of the third class (\emph{MR}~1962, p.~695); the postconciliar
calendar honours Raphael on 29 September. The collect asks for the
companion God gave Tobias:
\begin{quote}
\latin{Deus, qui beatum Raphaelem Archangelum Tobiae famulo tuo comitem
dedisti in via: concede nobis famulis tuis; ut ejusdem semper protegamur
custodia, et muniamur auxilio.} (\emph{MR}~1920, Mame 1922, p.~741)
\end{quote}
God gave the Archangel to his servant Tobias as a companion on the road,
and the Church asks that his servants be ever protected by the same
guardianship and strengthened by the same help. The Lesson is Raphael's
self-revelation to Tobias and his father: ``When thou didst pray with
tears, and didst bury the dead \dots\ I offered thy prayer to the Lord
\dots\ For I am the angel Raphael, one of the seven, who stand before the
Lord'' (Tob 12:12, 15; the Lesson is 12:7--15). The Gradual sings his victory over the demon,
\latin{Angelus Domini Raphael apprehendit et ligavit daemonem}, as
Tobias tells it: ``Then the angel Raphael took the devil, and bound him in
the desert of upper Egypt'' (Tob 8:3). The Gospel is the pool of the five
porches, where ``an angel of the Lord descended at certain times into the
pond and the water was moved. And he that went down first into the pond
after the motion of the water was made whole'' (John 5:1--4): the
medicine of God works through the angel of the pool. The Postcommunion
asks God to send Raphael to our help, \latin{et, quem tuae majestati
semper assistere credimus, tibi nostras exiguas preces benedicendas
assignet}, that the angel we believe to stand always before the divine
majesty may present our small prayers to God to be blessed (\emph{MR}~1920,
Mame 1922, p.~743). The Office of the day in Bute's Breviary reads a
sermon under Augustine's name on the book of Tobias, which draws the
moral of the story: ``Behold, dearly beloved brethren, how great is the
merit of almsgiving. Almsgiving hath earned an Angel for a servant''
(Bute, IV, p.~686).

\sectionguard
\subsection{The Faith the Liturgy Prays, by Degree}

The liturgy teaches the doctrine of the angels by praying it:

\begin{threestable}{Grade}{Proposition}{Liturgical witness}
Church teaching & The Church joins the angels' worship and venerates the
angels & Prefaces and Sanctus (Cummiskey 1861, pp.~xxviii--xxxii);
\emph{Directory} 215; CCC~335\\
Church teaching & The Church invokes the angels, Michael, Gabriel, and
Raphael by name, and the rest by order & Confiteor; Litany of the Saints
(\emph{MR}~1920, Holy Saturday); collects of 29 Sept., 24 Mar., 24 Oct.\\
Church teaching & God sends angels to guard each of the faithful, from
infancy to death & Collect of 2 Oct.; Compline collect; CCC~336\\
Common teaching & The angels present the prayers and the sacrifice of the
Church before God & Supplices (\emph{MR}~1920); incense prayer; Hilary,
lesson ix of 2 Oct.; \ST{III q.~83 a.~4 ad~9}\\
Common teaching & Michael defends the Church and receives the souls of the
just at death & Alleluia of 29 Sept.; Offertory of the Requiem; Lauds
antiphon ``Archangele Michael''\\
Common teaching & The angels assist the dying and lead the just into
paradise & \latin{Proficiscere}, \latin{Subvenite}, \latin{In paradisum}
(\emph{RR}~1872); \emph{Directory} 215\\
Common teaching & There are nine orders of angels & Gregory,
\emph{Hom.\ in Ev.}\ 34.7, lesson iv of 29 Sept.; Prefaces\\
Pious tradition & The Archangel appeared on Mount Gargano and asked that
God be worshipped there & Lessons of 8 May (\emph{BR}); Martyrology of
29 Sept.\ (Section~\ref{sec:devotion})\\
\end{threestable}
